\clearpage\fancyhead[L]{\small COMET.}\hypertarget{native-390-lower}{}\label{native-390-lower}
\begin{center}\textbf{COMET.}\par\vspace{4pt}\hrule\vspace{4pt}\fontsize{10}{12}\selectfont Göttingische gelehrte Anzeigen. No. 105. pp. 1041..1043. 3 July 1815.\par\vspace{4pt}\hrule\vspace{4pt}\end{center}

\par\indent Issue no. 55 of this journal [p. 382 of this volume] reported the first observations made at the local observatory of the comet now visible. We can now add the following two by Prof. \textsc{Gauss}, whose observations had been interrupted for rather a long time by circumstances that happened to intervene.

\begin{center}\fontsize{10}{12}\selectfont\setlength{\tabcolsep}{4pt}\renewcommand{\arraystretch}{1.15}\begin{tabular}{l|c|c|c}
1815 & Mean time & Right ascension & Declination.\\\hline
April 2 & $9^{\mathrm u}\;11^{\mathrm m}\;55^{\mathrm s}$ & $61^\circ\;2'\;27''$ & $45^\circ\;39'\;57''\;\mathrm{N.}$\\
June 12 & $10\;53\;31$ & $165\;3\;3$ & $52\;8\;25$\\
\end{tabular}\end{center}

\Needspace{23\baselineskip}\par\indent Of the numerous observations communicated by astronomers elsewhere, we give only those made by Mr. \textsc{Santini} at the Padua observatory. Ritter \textsc{Brera}, secretary of the local section of the Imperial and Royal Institute, sent them to Prof. \textsc{Gauss} in a letter dated 19 May.

\SourceLeaf{391}\begin{center}\fontsize{10}{12}\selectfont\setlength{\tabcolsep}{5pt}\renewcommand{\arraystretch}{1.15}\begin{tabular}{l|c|c|c}
1815 & \shortstack{Mean time\\in Padua} & \shortstack{Right\\ascension} & \shortstack{Northern\\declination}\\\hline
April 24 & $9^{\mathrm u}\;39^{\mathrm m}\;50^{\mathrm s}$ & $80^\circ\;49'\;33''$ & $56^\circ\;3'\;16''$\\
28 & $8\;57\;49$ & $85\;57\;54$ & $57\;37\;53$\\
May 1 & $8\;57\;22$ & $90\;16\;13$ & $58\;41\;7$\\
2 & $9\;8\;51$ & $91\;47\;15$ & $59\;0\;32$\\
2 & $9\;40\;22$ & $91\;48\;48$ & $59\;0\;40$\\
6 & $9\;29\;48$ & $98\;13\;9:$ & $\text{cloudy}$\\
8 & $9\;10\;55$ & $101\;40\;4$ & $60\;35\;36$\\
8 & $9\;45\;26$ & $101\;42\;13$ & $60\;37\;0$\\
10 & $9\;9\;11$ & $105\;14\;41$ & $60\;59\;45:$\\
11 & $8\;37\;5$ & $107\;4\;0$ & $61\;6\;4$\\
11 & $9\;22\;30$ & $107\;7\;18$ & $61\;7\;30$\\
12 & $9\;3\;20$ & $108\;57\;37$ & $61\;16\;32:$\\
12 & $9\;48\;15$ & $109\;2\;11$ & $61\;16\;7$\\
17 & $10\;15\;15$ & $118\;47\;45$ & $61\;30\;45$\\
17 & $10\;31\;5$ & $118\;49\;35$ & $61\;30\;47$\\
\end{tabular}\end{center}

\par\indent These observations were made with a ring micrometer on the \textsc{Adams} quadrant.

\par\indent Several astronomers have since determined new parabolic elements, all of which differ little from Prof. \textsc{Gauss}’s first elements, reported in the issue cited above. However well these elements fit the earlier observations, they all shared the same fate. The later observations diverged from them more quickly and more markedly than usual in calculations of this kind. This already pointed to a noticeable difference between the orbit and a parabola. As the comet is now growing fainter and observations are becoming increasingly difficult, Prof. \textsc{Gauss} thought it would be useful to make another correction to the parabolic elements to facilitate the observations still to be made. From the observations of 6 March and 25 April (both by Dr. \textsc{Olbers}) and that of 12 June, he obtained the following results:

\Needspace{12\baselineskip}\begin{center}\fontsize{10}{12}\selectfont Passage through perihelion, 25 April.\par $11^{\mathrm u}\;41^{\mathrm m}\;19^{\mathrm s}$ Mean time in Göttingen\end{center}
\begin{center}\fontsize{10}{12}\selectfont\setlength{\tabcolsep}{3pt}\renewcommand{\arraystretch}{1.15}\begin{tabular}{lcr}
Longitude of perihelion & $.\;.\;.\;.\;.\;.\;.\;.\;.\;.\;.\;.$ & $147^\circ\;35'\;55''$\\
Least distance & $.\;.\;.\;.\;.\;.\;.\;.\;.\;.\;.\;.\;.\;.$ & $1.23024$\\
Ascending node & $.\;.\;.\;.\;.\;.\;.\;.\;.\;.\;.\;.\;.$ & $82^\circ\;43'\;6''$\\
Inclination of orbit & $.\;.\;.\;.\;.\;.\;.\;.\;.\;.\;.\;.\;.$ & $44\;43\;13$\\
\end{tabular}\end{center}\begin{center}Prograde motion.\end{center}

\par\indent These elements give a very close fit to the three observations in all respects. Even so, they differ very considerably from the intervening observations (by \emph{as much as eleven minutes} in right ascension). This phenomenon can occur often and should come as no surprise. It follows, on the one hand, that the orbit departs considerably from a parabola, and on the other, that the three observations above are unsuitable for determining the conic section, since they leave it to some extent undetermined on their own. (The real reason is that around the time of the middle observation, the direction of the geocentric motion almost coincided with the direction of a great circle drawn from the comet toward the Sun. See the last articles in the first section of the second book of the \emph{Theor. Mot. C. C.}.) The following short ephemeris, calculated from these parabolic elements, can therefore be expected to depart noticeably from the actual motion right from the start. It will give right ascensions that are too small and declinations that are too large, but it will still be accurate enough for the limited purpose of making the comet easier to find:

\SourceLeaf{392}\begin{center}\fontsize{10}{12}\selectfont\setlength{\tabcolsep}{5pt}\renewcommand{\arraystretch}{1.15}\begin{tabular}{l|c|c}
\shortstack{$10^{\mathrm u}\;53^{\mathrm m}$\\in Göttingen} & \shortstack{Right\\ascension} & Declination\\\hline
June 20 & $174^\circ\;46'$ & $46^\circ\;47'\;\mathrm{N.}$\\
24 & $178\;55$ & $43\;56$\\
28 & $182\;41$ & $41\;3$\\
July 2 & $186\;7$ & $38\;8$\\
6 & $189\;16$ & $35\;16$\\
10 & $192\;11$ & $32\;25$\\
14 & $194\;53$ & $29\;40$\\
\end{tabular}\end{center}

\par\indent On 20 June the brightness is still $\frac{3}{10}$ of what it was on the day the comet was discovered, whereas on 14 July it is only half that brightness.

\par\indent Although it is still too early to determine the comet’s true orbit accurately, Prof. \textsc{Gauss} thought it would be of interest to make a preliminary determination. This would at least give a rough idea of the unexpectedly large departure of the orbit from a parabola and could also help in discussing and preparing the observations for a more precise determination. In keeping with that purpose, the calculation was done only roughly, based on just five isolated observations, with no account taken of precession, aberration, nutation or parallax. It would have been a waste of effort to make a scrupulous allowance for these effects in such a first determination.

\Needspace{13\baselineskip}\begin{center}
Preliminary \emph{elliptical} elements of the comet of 1815.
\end{center}

\begin{center}\fontsize{10}{12}\selectfont Passage through perihelion, 26 April $0^{\mathrm u}\;21^{\mathrm m}\;25^{\mathrm s}$\end{center}
\begin{center}\fontsize{10}{12}\selectfont\setlength{\tabcolsep}{3pt}\renewcommand{\arraystretch}{1.15}\begin{tabular}{lcr}
Longitude of perihelion & $.\;.\;.\;.\;.\;.\;.\;.\;.\;.\;.$ & $148^\circ\;58'\;48''$\\
Distance at perihelion & $.\;.\;.\;.\;.\;.\;.\;.\;.$ & $1.21349$\\
Ascending node & $.\;.\;.\;.\;.\;.\;.\;.\;.\;.\;.$ & $83\;26\;21$\\
Inclination of orbit & $.\;.\;.\;.\;.\;.\;.\;.\;.\;.\;.\;.\;.$ & $44\;30\;43$\\
Eccentricity & $.\;.\;.\;.\;.\;.\;.\;.\;.\;.\;.\;.\;.\;.$ & $0.933149$\\
\end{tabular}\end{center}

\par\indent This also gives a distance at aphelion of 35.091 and a period of $77\frac{1}{8}$ years, almost as long as that of \textsc{Halley}’s comet. These elements completely remove all those large discrepancies of the parabolic orbit, so Prof. \textsc{Gauss} believes he can already regard them as an approximation to the truth. Even if a more precise calculation later yields a noticeably longer period, it can hardly exceed 100 years. The comet can therefore already be regarded as one of the most remarkable phenomena in the annals of astronomy.
